% !TEX root = ../glossario.tex
% Lettera C del Glossario (vedi capitoli/a.tex per le istruzioni).
\phantomsection
\letteraglossario{C}

\termine{Cache}{Memoria temporanea ad alta velocità utilizzata per archiviare dati a cui si accede frequentemente, al fine di velocizzarne il recupero nelle operazioni successive.}

\termine{Capitolato}{Documento con cui un'azienda proponente descrive il prodotto che
  richiede nel progetto didattico: obiettivi, requisiti obbligatori e
  opzionali, vincoli, tecnologie suggerite e prodotti attesi.}

\termine{Ciclo di vita}{Insieme di fasi e attività che descrivono il processo di sviluppo di un prodotto software, dalla concezione alla dismissione.}

\termine{Code Coverage}{Percentuale del codice sorgente eseguita almeno una volta dai test automatici; indica quanta parte del codice è effettivamente messa alla prova.}

\termine{Code review}{Pratica in cui il codice sorgente scritto da uno o più sviluppatori viene esaminato sistematicamente da altri membri del team per individuarne difetti, migliorarne la qualità e assicurarne la conformità agli standard.}

\termine{Code Smell}{Caratteristica del codice che non è un errore ma segnala un possibile problema di progettazione o di manutenibilità, come metodi troppo lunghi o codice duplicato.}

\termine{Committente}{Organizzazione che commissiona un progetto software, definendo i requisiti e le specifiche di alto livello, e che ne verifica la realizzazione.}

\termine{Complessità ciclomatica}{Metrica che conta i cammini indipendenti nel flusso di controllo di un metodo o di una funzione: più è alta, più il codice è difficile da capire e da testare.}
